\subsection{Механические гироскопы}

Работа механических гироскопов базируется на законах механики вращательного движения твердого тела, в частности, на законе сохранения кинетического момента \cite{babaeva1973giroskopy}.

Гироскопический эффект проявляется в двух основных свойствах:

\begin{enumerate}
	\item Свойство инерциальности (устойчивости): Быстро вращающийся ротор (маховик) с большим кинетическим моментом стремится сохранить неизменным направление своей оси вращения в инерциальном пространстве.
	\item Свойство прецессии: При действии на гироскоп внешнего момента силы $\vec{M}$ его ось вращения совершает движение не в направлении действия момента, а в направлении, ему перпендикулярном. 
\end{enumerate}



\subsubsection{Конструктивные разновидности и особенности}

\begin{itemize}
	\item Гироскоп в кардановом подвесе (трехстепенной): Классическая конструкция, где ротор помещен в систему из двух или трех рам, изолирующих его ось от поворотов основания. Обеспечивает прямое измерение углов, но имеет существенные недостатки: трение в осях подвеса, ограниченный угол поворота рам, сложность конструкции.
	
	\item Динамически настраиваемый гироскоп (ДНГ): Ротор вращается внутри вращающегося карданова подвеса. При определенном соотношении скоростей («условии настройки») момент сил инерции, действующих на рамки, уравновешивается, и гироскоп становится астатическим. ДНГ не имеет ограничений по углу поворота, но требует сложной системы управления.
	
	\item Поплавковый интегрирующий гироскоп: Чувствительный элемент (ротор в подвесе) помещен в герметичный поплавок, который плавает в высокоплотной жидкости. Это обеспечивает гидростатическую разгрузку осей, демпфирование колебаний и термостабилизацию.
\end{itemize}

Исторически механические гироскопы являются основой инерциальной навигационной системы в авиации, ракетостроении и космонавтике. В настоящее время практически полностью вытеснены оптическими системами, за исключением некоторых специальных применений (например, гиростабилизаторы).

К достоинствам можно отнести высокую теоретическую и практическую точность у поплавковых моделей, прямое измерение углов ориентации и высокую надёжность в условиях специфических воздействий таким как перегрузки, радиация и электромагнитные импульсы.

Из недостатков можно упомянуть дрейф гироскопа из-за трения подвижных частей, большие габариты и вес готовой конструкции, а также ограничения по углам поворота гироскопа.

